\documentclass{article}
\usepackage{hyperref}
\usepackage{datetime}
\usepackage{graphicx}
\usepackage{listings}
\usepackage{color}
\usepackage{todonotes}
\usepackage{lmodern}
\usepackage[T1]{fontenc}
\hypersetup{
    colorlinks=true,
    linktoc=all,
    linkcolor=blue,
}
\lstset{frame=tb,
  language=Java,
  aboveskip=3mm,
  belowskip=3mm,
  showstringspaces=false,
  columns=flexible,
  basicstyle={\small\ttfamily},
  numbers=none,
  numberstyle=\tiny\color{gray},
  keywordstyle=\color{blue},
  commentstyle=\color{dkgreen},
  stringstyle=\color{mauve},
  breaklines=true,
  breakatwhitespace=true,
  tabsize=3
}
\author{Peter Mezei}
\title{Gardenzilla \\ development notes}
\pagestyle{myheadings}
\markright{Gardenzilla\hfill Development Notes \hfill}
\begin{document}
\maketitle
\newpage
\thispagestyle{empty}
\begin{figure}
  \centering
  \includegraphics[width=3cm]{logo.png}
\end{figure}
\begin{center}
  {\textit  Gardenzilla development notes, questions and future development directions. The main goal of this document is to document the development steps and collect all the ideas for future evaluation.

    \vspace{1cm} All right reserved \the\year{}

    All the logos, images and all the other\\ used visual materials are copyrighted.
  }
\end{center}
\newpage
\thispagestyle{empty}
\tableofcontents
\newpage
\section{Introduction}
\newpage

% section MODULES
% ***************
\section{Modules}

\subsection{Customers}

We manage customers as independent data objects, and manage a customer
object through its implemented methods, or via public functions inside
core::customer module.

\begin{figure}[!htb]
  \caption{Customer object}
  \begin{verbatim}
    pub struct Customer {
      id: String,
      related_users: Vec<String>,
      name: String,
      tax_number: String,
      address: InvoiceAddress,
      phone: String,
      email: String,
      date_created: DateTime<Utc>,
      created_by: String,
    }
  \end{verbatim}
\end{figure}

Where:

\begin{description}
  \item[id] \hfill
    \\ Customer unique identifier \todo {Describe and set the way we manage IDs.}
  \item[related\_users] \hfill
    \\ vector of related user ID. User and Customer objects are independent
    of each other, as we can have customer without a user, and we can have
    - at least in theory - a user without a customer. Customer can
    purchase, User can login and interact withing the client. To manage
    Customer <-> User relations, we store them as a related ID list in
    each other. Here in a Customer object thats why we have this
    related\_users vector.

  \item[name] \hfill
    \\ name as string of the customer

  \item[tax\_number] \hfill
    \\ tax number as string
  \item[address] \hfill
    \\ address as InvoiceAddress object \todo{display InoviceAddress}
  \item[phone] \hfill
    \\ phone number as string
  \item[email] \hfill
    \\ email as string
  \item[date\_created] \hfill
    \\ date as chrono::DateTime<Utc>
  \item[created\_by] \hfill
    \\ userid as string, who created this object \todo{describe datecreated
      and created by somewhere as we use them everywhere}
\end{description}

\subsubsection{Methods}

\begin{description}
  \item[new(\dots) -> Self] \hfill
    \\ creates a new instance of Customer. ID cannot be changed after
    creation
  \item[get\_id() -> \&str] \hfill
    \\ returns customer id as \&str
  \item[get\_name() -> \&str] \hfill
    \\ returns customer name as \&str
  \item[set\_name(name: String)] \hfill
    \\ set customer name to the given new one
  \item[has\_user() -> bool] \hfill
    \\ returns bool if there is at least one related user
  \item[get\_users() -> \&Vec<String>] \hfill
    \\ returns related user vector reference
  \item[remove\_user(userid: \&str)] \todo{should return bool, or result} \hfill
    \\ remove a user by its userid
  \item[get\_tax\_number() -> \&str] \hfill
    \\ returns tax number as \&str
  \item[set\_tax\_number(tax\_number: String)] \hfill
    \\ set tax number
\end{description}

\newpage
\todototoc
\listoftodos

\newpage
\thispagestyle{empty}
\vspace*{5cm}
\begin{center}
  {\textit Gardenzilla\\ Development Notes \\ \vspace{5mm} All right reserved \\ Peter Mezei \\\the\year{}}
\end{center}
\end{document}
